% Analyse qualité et efficacité du dépôt VetNutriMP (Compose Multiplatform).
% Compilation : pdflatex analyse-qualite-efficacite-2026-10-25.tex (deux passes), paquets TeX Live
% latex-recommended, latex-extra, lang-french.
\documentclass[11pt,a4paper]{article}
\usepackage[T1]{fontenc}
\usepackage[utf8]{inputenc}
\usepackage{lmodern}
\usepackage[french]{babel}
\usepackage[margin=2.2cm]{geometry}
\usepackage{booktabs}
\usepackage{longtable}
\usepackage{tabularx}
\usepackage{array}
\usepackage{enumitem}
\usepackage{xcolor}
\usepackage[colorlinks=true,linkcolor=blue!50!black,urlcolor=blue!50!black]{hyperref}

\setlist{nosep}
\sloppy
\setlength{\emergencystretch}{3em}
\newcommand{\fichier}[1]{\texttt{\small #1}}
\newcommand{\code}[1]{\texttt{#1}}
\newcommand{\sev}[1]{\textbf{#1}}
\newcolumntype{Y}{>{\raggedright\arraybackslash}X}

\title{VetNutriMP : analyse qualité et efficacité du projet\\
  \large Version 3.3.18 (318), schéma Room 41, Kotlin 2.2.20, Compose Multiplatform 1.7.0}
\author{Sébastien Lefebvre}
\date{25 octobre 2026 --- commit \code{24b85e7}}

\begin{document}
\maketitle

\section*{Synthèse}

Le projet est fonctionnellement riche (quatre cibles, 29 tables, 24 migrations, exploration
multiration, mode troupeau, export PDF, partage en ligne) et plusieurs choix structurants sont
sains : aucune migration destructive, base mise de côté plutôt qu'effacée en cas d'échec,
sauvegarde atomique par fichier temporaire puis renommage, chaîne de migrations testée sur une
vraie base SQLite, 789 tests unitaires, intégration continue qui conditionne l'empaquetage aux
tests et contrôle la dérive des schémas Room exportés, aucune clé en dur dans les sources.

Le revers est net. Les règles d'architecture écrites dans \fichier{CLAUDE.md} ne sont pas
appliquées : 25 fichiers de vues reçoivent des repositories en paramètre, le mécanisme
\code{AppNavRepositories} officialise ce contournement, et des règles métier (arrondi des
quantités) vivent dans un fichier de dialogue importé par la couche \code{Data}. Sept fichiers
dépassent 2\,500 lignes ; \code{SettingsView} est un composable de 2\,546 lignes. Les accès base
sont dominés par des schémas N+1 (références, animaux, consultations) et par le chargement du
catalogue entier (environ 10\,000 aliments avec leurs nutriments) pour obtenir un simple compteur,
le tout relancé toutes les dix minutes par la sauvegarde automatique.

Quatre défauts méritent une correction immédiate : la suppression d'un aliment utilisé dans une
ration échoue sur une contrainte \code{NOT NULL} et laisse l'aliment sans ses nutriments ;
\code{NumberUtils.format} produit des valeurs fausses dans l'export PDF (2,3 affiché « 2.2 »,
négatifs illisibles) ; deux fichiers JSON à la racine du dépôt contiennent 132 noms de
propriétaires réels ; le cache global d'évaluation d'équations n'est pas synchronisé et sa clé
repose sur un \code{hashCode} sujet aux collisions.

Enfin, le dépôt lui-même pèse : deux AAB de release (40 Mo) commités malgré le \fichier{.gitignore},
cinq copies d'un JSON de 20,7 Mo, 52 Mo de jeux de données bruts à la racine, sept fichiers
Kotlin orphelins, deux répertoires \fichier{commonMain/} et \fichier{iosMain/} hors de tout source
set, et 18 fichiers Markdown à la racine dont plusieurs décrivent des optimisations jamais
réalisées.

\section{Méthode et périmètre}

Lecture statique du dépôt \code{vinlmoe/VetNutriMP} à la révision \code{24b85e7}, sans
modification du code. Quatre revues indépendantes ont été menées puis recoupées : architecture et
maintenabilité, performance et efficacité, robustesse, données et sécurité, tests, build, CI et
hygiène du dépôt. Chaque constat retenu cite un fichier et une ligne ; les trois constats les plus
lourds (format numérique, contrainte de clé étrangère, données personnelles) ont été revérifiés
indépendamment.

Limite : la suite de tests n'a pas pu être exécutée dans l'environnement d'analyse, dont la
politique réseau bloque le dépôt Maven de Google nécessaire au plugin Android. Les affirmations
sur les tests reposent donc sur leur lecture et sur la configuration CI, pas sur une exécution.

Le dépôt frère ThesisManager dispose de son propre rapport daté du même jour et n'est pas traité
ici.

\section{Chiffres clés}

\begin{center}
\begin{tabular}{lr}
\toprule
Lignes Kotlin (\fichier{composeApp/src}) & 136\,653 \\
Fichiers Kotlin \code{commonMain} : View / Data / Utils / ViewModel / Repository & 101 / 56 / 44 / 26 / 19 \\
Fonctions \code{@Composable} & 392 \\
Composables de plus de 300 lignes & 46 sur 381 \\
Fichiers de plus de 2\,000 lignes & 10 \\
Tests (\code{@Test}) : commonTest / desktopTest / iosTest & 723 / 60 / 6 \\
ViewModels couverts par un test & 2 sur 26 \\
Tests d'interface Compose & 0 \\
Tables Room / migrations / version de schéma & 29 / 24 / 41 \\
\code{catch} vides en \code{commonMain} & 67 \\
Opérateurs \code{!!} en \code{commonMain} & 142 \\
\code{LazyColumn}/\code{LazyRow} sans \code{key} & 48 sur 65 \\
Fichiers Markdown à la racine & 18 \\
Taille du clone (partiel) & 59 Mo \\
Commits & 103 \\
\bottomrule
\end{tabular}
\end{center}

Les compteurs de \fichier{CLAUDE.md} sont périmés : 51 vues annoncées pour 101 réelles, 23
ViewModels pour 26, 23 écrans pour 26, JSON initial « ~11 Mo » pour 20,7 Mo, Room « 2.7.0 » pour
\code{2.7.0-rc01}.

\section{Architecture et maintenabilité}\label{sec:architecture}

\subsection*{Ce qui est fait correctement}

\begin{itemize}
\item Aucun \code{MutableStateFlow} n'est exposé publiquement : tous passent par \code{asStateFlow()}.
\item Les 15 fichiers par plateforme dans \fichier{Utils/} sont bien du \code{expect}/\code{actual}
  (36 \code{expect}, 59 \code{actual}), pas du copier-coller, à deux exceptions près.
\item La règle « \code{NutrientResolver} seul point d'entrée » est respectée pour les nutriments.
\item Les moteurs purs \fichier{Data/Troupeau.kt} et \fichier{Data/ExplorationMultiration.kt} sont
  séparés de l'interface et testés (27 et 17 tests).
\end{itemize}

\subsection*{Constats}

\begin{description}[style=nextline]
\item[\sev{Critique} --- Les vues accèdent directement aux repositories.]
  69 paramètres de type \code{XxxRepository} sur 25 fichiers de \fichier{View/} et
  \fichier{Components/}, 34 appels effectifs. \code{AppNavRepositories}
  (\fichier{Navigation/AppNavModels.kt:50-58}) injecte six repositories dans sept écrans.
  \fichier{View/StartupScreen.kt:180-194} lit et réécrit des référentiels dans un
  \code{LaunchedEffect}. \fichier{View/AnimalDetailView.kt:433-453} construit un
  \code{ExportImportRepository} dans la vue et contient 13 \code{withContext(AppDispatchers.IO)}.
  \fichier{View/SettingsView.kt:2702} instancie un \code{PreferencesRepository} dans un
  \code{remember}. Cause racine : les ViewModels exposent leurs repositories en \code{internal val}
  (\fichier{ViewModel/SettingsViewModel.kt:33-46}, \fichier{AnimalListViewModel.kt:45-51}).
  \emph{Remède} : rendre ces propriétés privées et supprimer \code{AppNavRepositories} ; le
  compilateur listera les sites à migrer. Les ViewModels \code{ConseilsViewModel} et
  \code{StartupViewModel} existent déjà et ne sont pas instanciés.

\item[\sev{Majeur} --- Logique métier dans les vues, dépendance inversée Data $\rightarrow$ View.]
  \code{arrondirQuantiteSelonRegles} (règle d'arrondi des quantités de ration) est définie dans
  \fichier{View/AnalNut/MultiNutrientAdjustmentDialog.kt:2211} et importée par
  \fichier{Data/ExplorationMultiration.kt:13} et \fichier{Data/RationConstraintAdjuster.kt:17}.
  \fichier{View/AnalyseGraphique/AnalyseGraphiqueUtils.kt} et
  \fichier{View/AnalyseGraphiqueAlimentsView.kt:87-500} contiennent des calculs nutritionnels
  suspendus prenant des repositories. \emph{Remède} : déplacer vers \fichier{Data/} et tester.

\item[\sev{Majeur} --- Trois implémentations du poids métabolique.]
  \fichier{Data/CalculMetabolique.kt:46} applique l'équation BW du référentiel ;
  \fichier{Data/RationCalculator.kt:52-57} et
  \fichier{ViewModel/CrossConsultationAnalysisViewModel.kt:409} codent \code{pow(0.75)} en dur.
  Résultats différents selon l'écran pour une espèce à exposant 0,67.
  \code{CalculMetabolique.kt} n'a par ailleurs aucun test.

\item[\sev{Majeur} --- Fichiers monolithiques.]
  \code{SettingsView} est un composable de 2\,546 lignes dont l'onglet Importation occupe
  2\,300 lignes en ligne, alors que \fichier{SettingsSections/} existe. \code{RationsView}
  (1\,901 lignes, 16 booléens \code{showXxxDialog}), \code{StartupScreen} (1\,724 lignes),
  \code{AnimalDetailView} (deux layouts large et étroit de 1\,070 et 1\,169 lignes qui dupliquent
  la même composition), \code{ImportUtils} (objet de 50 méthodes mêlant quatre domaines et du code
  de test), \code{AnimalDetailViewModel} (33 \code{StateFlow}, 88 fonctions). Fonctions non
  composables de plus de 300 lignes : \code{ExportImportRepository.importAll} (648),
  \code{DatabaseAnimalRepository.importAnimals} (465), \code{DatabaseFoodRepository.importFoods} (400).

\item[\sev{Majeur} --- Cycle de vie des ViewModels incohérent.]
  Huit classes étendent \code{androidx.lifecycle.ViewModel} mais sont créées par \code{remember}
  dans \fichier{App.kt:108-258} sans \code{ViewModelStore} : \code{onCleared()} n'est jamais
  appelé. Treize autres créent un \code{CoroutineScope} maison ; cinq n'ont aucune méthode
  d'annulation. \fichier{ViewModel/RecipeEditViewModel.kt:30} mute de l'état Compose
  (\code{mutableStateOf}) depuis un scope \code{IO}. \code{EquationViewModel} reçoit un DAO
  directement (\fichier{App.kt:166}).

\item[\sev{Majeur} --- Règle \code{withContext(AppDispatchers.IO)} non suivie.]
  Cinq repositories entiers l'ignorent : \code{ConseilRepository} (15 fonctions),
  \code{HtmlSectionRepository} (10), \code{DatabaseEquationRepository} (6),
  \code{ExamGradingRepository} (5), \code{PreferencesRepository} (5) ; \code{ExportImportRepository}
  encode 20 Mo de JSON hors IO. Room exécute ses DAO hors thread principal, le risque réel est
  faible, mais la règle doit être appliquée ou reformulée.

\item[\sev{Majeur} --- Navigation sans pile.]
  \fichier{Navigation/AppNavController.kt} porte un seul \code{screen} et onze identifiants
  sélectionnés ; 65 lambdas de retour codent la destination en dur dans \fichier{AppNavHost.kt}.
  \code{detailReturnScreen} est un contournement ponctuel. \code{Screen.TestYellowBox} est livré
  en production.

\item[\sev{Mineur} --- Code mort.]
  \fichier{Example/} et \fichier{Examples/} (5 fichiers), \fichier{Utils/ImportTester.kt},
  \fichier{Repository/TestApiEnvelope.kt} (vide), \fichier{Repository/ExportImportRepository.kt.disabled}
  (59 Ko), \fichier{Repository/ConseilRepositoryExample.kt}, \code{Greeting} et
  \fichier{View/MainViewModel.kt} (gabarit KMP), \fichier{JsonFileImport.kt}, cinq ViewModels
  jamais instanciés, \fichier{Enumer/OptimizedEspece.kt} sans appelant, \code{LruCache} (thread-safe,
  TTL) utilisé uniquement par ses tests. Environ 1\,000 lignes de bruit.

\item[\sev{Mineur} --- Duplications.]
  \code{arrondirPourAffichage}/\code{calculateAdaptiveRange} copiés verbatim entre
  \fichier{AnalyseGraphiqueUtils.kt:247-280} et \fichier{AnalyseGraphiqueAlimentsView.kt:516-598} ;
  deux \code{formatDecimal} aux sémantiques différentes ; résolution d'espèce depuis un label
  dupliquée en cinq endroits ; \code{PlatformCrypto.desktop.kt} identique à
  \code{PlatformCrypto.android.kt} (un source set \code{jvmMain} intermédiaire manque).
\end{description}

\section{Efficacité et performance}\label{sec:performance}

\subsection*{Ce qui est fait correctement}

\begin{itemize}
\item Caches d'aliments protégés par \code{Mutex} et invalidés à chaque écriture
  (\fichier{Repository/DatabaseFoodRepository.kt:74-107}).
\item Import initial idempotent : seule la version du JSON est lue (4 Ko) avant de décider
  d'importer (\fichier{ViewModel/SettingsViewModel.kt:598-623}).
\item Export JSON en flux sur desktop et Android ; Ktor avec délais et reprise exponentielle sur
  les trois plateformes ; R8 activé ; debounce de 300 ms sur la recherche d'aliments.
\end{itemize}

\subsection*{Constats}

\begin{description}[style=nextline]
\item[\sev{Critique} --- \code{getAllReferenceEv()} : N+1 à trois niveaux.]
  \fichier{Repository/DatabaseReferenceEvRepository.kt:186-220, 585-720} : par référence, une
  requête par équation, par coefficient et par nutriment (puis par biblio). Plusieurs milliers de
  requêtes par appel ; appelé 14 fois dont cinq uniquement pour \code{.size}
  (\fichier{View/StartupScreen.kt:146, 418, 1388, 1500}, \fichier{SettingsViewModel.kt:589}).
  \emph{Remède} : quatre requêtes \code{IN (:ids)} et \code{groupBy}, plus un \code{COUNT(*)}.

\item[\sev{Critique} --- Catalogue entier chargé pour un compteur ou des UUID.]
  \code{getAllFoods().size} juste après un import qui vient de vider le cache
  (\fichier{StartupScreen.kt:145, 415, 1385, 1499}), \code{getAllFoods().map \{ it.uuid \}}
  (\fichier{ExportImportRepository.kt:951, 1062}) alors que \code{getAllFoodIds()} et
  \code{getFoodsCount()} existent. Chaque appel reconstruit environ 10\,000 \code{AlimentEv} avec
  leurs valeurs nutritionnelles.

\item[\sev{Majeur} --- Sauvegarde automatique toutes les dix minutes sans détection de changement.]
  \fichier{Service/BackupService.kt:34, 65-74, 149} exporte toute la base : \code{getAllAnimals()}
  (N+1), consultations (N+1 imbriqué), \code{forceRefresh()} qui annule le cache de cinq minutes
  puis \code{getAllFoods()}, \code{getAllReferenceEv()} (N+1), sérialisation d'environ 20 Mo. Sur
  iOS la chaîne JSON complète est construite en mémoire (\fichier{JsonStreamUtils.ios.kt:10}).
  \emph{Remède} : indicateur de modification via \code{DatabaseChangeNotifier}, catalogue
  d'origine exclu ou sauvegardé séparément.

\item[\sev{Majeur} --- Autres N+1.]
  Poids par animal (\fichier{Repository/DatabaseAnimalRepository.kt:55-67}),
  \code{getRacesBySpecies} qui charge tous les animaux (\code{:647-655}), variables et rations par
  consultation puis aliments par ration (\fichier{Repository/ConsultationRepository.kt:102-118}),
  \code{getFoodById} dans une triple boucle à l'import (\code{:540, 599}).

\item[\sev{Majeur} --- Import d'aliments sans transaction englobante.]
  \fichier{DatabaseFoodRepository.kt:1083-1140} : insertions par lots mais aucun
  \code{useWriterConnection}/\code{immediateTransaction} ; suppression de l'énergie par espèce
  une requête par aliment (\code{:1123-1131}) ; repli à une transaction par valeur nutritionnelle
  (\code{:1113-1120}).

\item[\sev{Majeur} --- Recherche \code{LIKE '\%…\%'} non indexable.]
  \fichier{DataBase/CommonDao.kt:336-352} balaye la table FOOD à chaque recherche ; les index sur
  \code{name}, \code{brand}, \code{especesJson} ne servent pas ; \code{especesJson LIKE} alors
  que la table de jonction \code{ESPECES\_ALIMENTS} existe. Le \code{searchCache} retient jusqu'à
  50 listes de 5\,000 \code{AlimentEv}. Index \code{uuid} redondants avec la clé primaire
  (\fichier{Entity.kt:20, 139, 156}).

\item[\sev{Majeur} --- Travail lourd sur le thread principal.]
  \code{AnimalDetailViewModel} tourne sur \code{Main} (\code{:87}) et appelle
  \code{analyserRation} sans \code{withContext(Default)} aux lignes 771 et 1570 (contrairement
  aux lignes 1987 et 2623). À chaque émission du catalogue, copie de 10\,000 \code{AlimentEv}
  (\code{:306-316}) et indexation normalisée dans \code{FoodListViewModel} (\code{:92-100}).
  \code{Espece.getFromString} recompile une \code{Regex} et reconstruit deux maps à chaque appel
  (\fichier{Enumer/Espece.kt:47-51}), soit environ 20\,000 fois par rafraîchissement.
  Sur Android, l'ouverture Room avec migrations est bloquante dans \code{onCreate}
  (\fichier{AppDatabase.kt:181}, \fichier{MainActivity.kt:39}).

\item[\sev{Majeur} --- Filtrage de 10\,000 aliments pendant la composition.]
  \fichier{View/components/FoodSearchComponent.kt:148-368} filtre et trie dans un
  \code{remember} à 14 clés, à chaque frappe sans debounce ; \code{:102} alloue
  \code{foods.map \{ it.uuid \}} à chaque recomposition comme clé de \code{LaunchedEffect} ;
  \code{:918} \code{items(foods)} sans \code{key}. 48 listes paresseuses sur 65 n'ont pas de clé.
  223 \code{collectAsState()} et aucun \code{collectAsStateWithLifecycle} malgré la dépendance.

\item[\sev{Majeur} --- Cache d'équations non synchronisé et clé faillible.]
  \fichier{Utils/EquationEvaluator.kt:36, 171, 199} : \code{LruMap} documentée non thread-safe,
  accédée depuis \code{Default} et \code{Main} ; clé \code{"\$expression:\$\{variables.hashCode()\}"},
  une collision renvoie un résultat faux. \code{NutrientResolver.resolvedCache} est un
  \code{HashMap} global non borné, écrit depuis IO ; \code{labelToNutrient by lazy} fige les
  nutriments personnalisés enregistrés ensuite.

\item[\sev{Majeur} --- JSON initial de 20,7 Mo embarqué deux fois dans l'APK.]
  \fichier{build.gradle.kts:204} ajoute \fichier{commonMain/resources} aux ressources Android
  alors qu'Android lit uniquement l'asset. Assets jamais référencés : \fichier{vetfood.json}
  (1,7 Mo), \fichier{references.json} (1 Mo). Sur iOS, parsing par chaîne complète
  (\fichier{ResourceReader.ios.kt:64-70}), soit environ 40 Mo de pic.

\item[\sev{Mineur} --- Divers.]
  Cache \code{rationAnalysisCache} de fait mort (le chemin utilisé ne passe pas par lui, et
  \code{analyserRation(ration, consultation)} n'a aucun appelant). Deux
  \code{DatabaseFoodRepository} instanciés au démarrage desktop (\fichier{main.kt:103-111} et
  \fichier{AppContainer.kt:55-62}), chacun lisant la table FOOD. \code{HttpClient} créé à chaque
  partage et jamais fermé (cinq sites). Dépendances \code{androidx.paging} déclarées et jamais
  utilisées. 14 \code{println} en \code{commonMain}. Scope de \code{DatabaseFoodRepository} sans
  \code{SupervisorJob} (\code{:67}).
\end{description}

\subsection*{Documents d'optimisation contredits par le code}

\fichier{RESUME\_OPTIMISATIONS\_PRIORITAIRES.md} affirme « État actuel : implémenté » pour
\code{OptimizedNutrientResolver.kt}, \code{OptimizedResourceReader.kt}, \code{CacheFactory},
trois scripts shell et des maps O(1) dans \code{UnitEnum} : aucun n'existe (\code{UnitEnum.byId}
reste une boucle). \fichier{OPTIMISATION\_MEMOIRE.md} décrit \code{-Xmx16g} et des options
\code{-Xno-optimize} absentes de \fichier{gradle.properties} (\code{-Xmx8g}).
\fichier{ANALYSE\_OPTIMISATION\_MEMOIRE.md} propose du SQL PostgreSQL (\code{PARTITION},
\code{CREATE INDEX CONCURRENTLY}) impossible en SQLite. \fichier{PERFORMANCE\_IMPROVEMENTS.md}
annonce « 0 println » et des clés sur les listes : partiellement vrai (RationsView et
AnimalDetailView seulement). Ces documents doivent être archivés ou réécrits.

\section{Robustesse, données et sécurité}\label{sec:robustesse}

\subsection*{Ce qui est fait correctement}

\begin{itemize}
\item Aucun \code{fallbackToDestructiveMigration} ; chaîne 17 $\rightarrow$ 41 continue ; échec
  de migration $\Rightarrow$ base mise de côté (\fichier{AppDatabase.kt:183-199}) ; \code{.bak} à
  chaque démarrage, rotation des bases corrompues, parking d'une base plus récente que
  l'application (\fichier{DatabaseGuard.kt}).
\item Sauvegarde : fichier temporaire puis renommage, rotation sur dix fichiers, copie long terme,
  \code{Mutex}, sauvegarde de secours avant restauration, confirmation utilisateur.
\item Aucune clé en dur dans l'arbre de travail ; clés jsonbin injectées à la compilation depuis
  \fichier{local.properties} ou l'environnement ; \fichier{local.properties} ignoré par Git.
\item Évaluation d'équations : division par zéro levée, NaN et infinis filtrés ; moteurs troupeau
  et multiration protégés contre les effectifs et densités nuls.
\end{itemize}

\subsection*{Constats}

\begin{description}[style=nextline]
\item[\sev{Critique} --- Suppression d'un aliment utilisé : échec et perte partielle.]
  \fichier{DataBase/Entity.kt:255-270} : \code{ALIMENTS.refAlimUnif} est \code{NOT NULL}
  (\code{String = ""}) mais la clé étrangère vers FOOD est \code{ON DELETE SET NULL}. SQLite lève
  \code{NOT NULL constraint failed}. \code{DatabaseFoodRepository.deleteFood} (\code{:1537-1547})
  supprime d'abord nutriments et indications hors transaction, puis échoue : l'aliment reste sans
  nutriments. \code{FoodListViewModel.deleteFood} (\code{:278-284}) n'a pas de \code{try/catch}
  dans son \code{launch} : exception non gérée sur Android.
  \emph{Remède} : colonne nullable (migration 42) ou \code{RESTRICT} avec message ;
  \code{@Transaction} ; erreur remontée à l'interface.

\item[\sev{Critique} --- \code{NumberUtils.format} est faux.]
  \fichier{Utils/NumberUtils.kt:6-22} tronque l'erreur flottante :
  \code{format(2.3, 1) = "2.2"}, \code{format(4.35, 2) = "4.34"},
  \code{format(-1.25, 2) = "-1.-25"}, NaN $\Rightarrow$ « 0 ». Vérifié numériquement. 28 usages
  dont \fichier{Export/HtmlDocumentBuilder.kt:119-124} (export PDF des rations) et
  \fichier{View/ConsultationDetailView.kt:368-444} (poids, kcal/j).
  \emph{Remède} : remplacer par \code{TextUtils.formatDecimal}, qui est correct.

\item[\sev{Critique} --- Données personnelles dans le dépôt.]
  \fichier{ani.json} (18 Mo) et \fichier{animaux.json} (19 Mo) à la racine contiennent 223 animaux
  et 132 noms distincts de propriétaires (\code{nomProprio}). Dépôt poussé sur GitHub.
  \emph{Remède} : retirer, purger l'historique, anonymiser (\code{AnonymizationUtils} existe).

\item[\sev{Majeur} --- Restauration et import non transactionnels, et trompeurs.]
  \code{importAll} (\fichier{ExportImportRepository.kt:509-1160}) insère unitairement sans
  transaction ; la section Équations (\code{:603-613}) n'a aucun \code{try/catch} et interrompt
  l'import à mi-chemin ; \code{envelope.version} n'est jamais vérifié ; le dialogue annonce
  « remplacera toutes les données » (\fichier{View/BackupRestoreView.kt:203}) alors qu'il s'agit
  d'une fusion sans suppression (\code{:964-970}). À l'import d'aliments, les échecs d'insertion
  des nutriments sont avalés (\fichier{DatabaseFoodRepository.kt:427, 444, 458, 1119, 1134}) et
  l'aliment compté comme importé.

\item[\sev{Majeur} --- Erreurs silencieuses sur des enregistrements cliniques.]
  Dans \code{AnimalDetailViewModel}, neuf fonctions de modification de consultation (référentiel,
  maladies, coefficients, variables ; lignes 1685 à 2334) entourent \code{updateConsultation}
  d'un \code{catch} vide, et \code{updateConsultation} lui-même ne fait que \code{printStackTrace}
  (\code{:1078-1080}) : l'interface affiche une valeur non enregistrée sans retour.
  \code{applyRecipeToRation} (\fichier{Repository/ConsultationRepository.kt:30-35}) applique une
  recette partiellement sans erreur. 67 \code{catch} vides au total.

\item[\sev{Majeur} --- Clés jsonbin embarquées dans les binaires.]
  \code{internal const val} générées dans \code{AppSecretsGenerated.kt} : extractibles par
  \code{strings} du JAR, de l'APK ou de l'IPA. Elles permettent de créer des bins et de lire tout
  bin privé du compte (partage de dossiers entre confrères). Chiffrement AES-CBC sans
  authentification (\fichier{Utils/CryptoUtils.kt:19}). \emph{Remède} : relais serveur, ou
  chiffrement obligatoire en AES-GCM avec bins éphémères. Android : \code{allowBackup="true"}
  sans règles d'exclusion, la base patients part dans la sauvegarde cloud.

\item[\sev{Majeur} --- Migrations.]
  Aucune migration avant la version 17 ; migration 23 $\rightarrow$ 24 avec
  \code{DROP TABLE IF EXISTS HTML\_SECTIONS} (\fichier{AppDatabase.kt:449-455}) ;
  \code{RATIONS.refRationParente} sans clé étrangère (étapes de plan évolutif orphelines) ; le test
  de chaîne annonce « 17 $\rightarrow$ 36 » dans son commentaire. Sur Android,
  \code{MainActivity.onCreate} rappelle \code{getRoomDatabase} à chaque recréation sans fermer
  l'instance précédente.

\item[\sev{Majeur} --- \code{!!} sur des entrées externes.]
  \fichier{Utils/ImportUtils.kt:582, 601, 619} (\code{value["value"]!!} sur JSON importé),
  \fichier{Utils/LegacyV2Mappings.kt:151-165} (\code{find(line)!!} sur scripts hérités),
  \fichier{Data/RationAnalyzer.kt:124-132} (clé de ratio codée par libellé).

\item[\sev{Mineur} --- Divers.]
  Variable d'environnement \code{VETNUTRI\_TEST\_DATABASE\_DIR} honorée en production
  (\fichier{desktopMain/DataBase/Database.kt:8}) ; \code{createDatabaseBackup()} stub retournant
  \code{true} ; \code{Json.isLenient} et \code{coerceInputValues} sur les imports utilisateur ;
  \code{moveFile} sans \code{ATOMIC\_MOVE} ; échec d'écriture de \code{startup\_complete} avalé
  (écran de démarrage réaffiché en boucle) ; \fichier{GUIDE\_COMPATIBILITE\_16KB.md} cite NDK 27
  et un script inexistant alors que le build utilise NDK 28 ; exclusion de
  \code{libimage\_processing\_util\_jni.so} à valider sur appareil.
\end{description}

\section{Tests}

\begin{center}
\begin{tabularx}{\textwidth}{lYY}
\toprule
Couche & Couvert & Non couvert \\
\midrule
ViewModel (26) & \code{AnimalDetailViewModel} (plan évolutif, concurrence), 3 tests sur une fonction pure de \code{CrossConsultation} & 24 autres, dont MultiRation, FoodEdit, Import, Settings, Startup, Backup \\
Repository (19) & InMemory Animal/Equation, \code{DatabaseFoodRepository} (vrai Room), Consultation évolutive, ExportImport (4 fichiers), recettes & \code{DatabaseAnimalRepository}, Conseil, HtmlSection, ExamGrading, Preferences \\
Service (9) & Excel aliments (22), Backup (11), File partiel, migration V2 (1) & \code{StartupService}, CrossRef, CSV notes, JsonShare \\
Data (moteurs) & Troupeau (27 + 15 JSON), Multiration (17), RationAnalyzer (11), RationCalculator (8), PlanEvolutif (9) & \code{CalculMetabolique} (BEE, poids métabolique), RequirementAnalyzer \\
DataBase & Chaîne de migrations sur base v17 réelle, persistance, garde, rotation (26 tests) & Migrations 37 à 41 à confirmer dans le test de chaîne \\
Export & Prescription plan évolutif (2) & \code{PdfExporter} \\
View (101) & 8 fichiers testant des fonctions pures placées dans \fichier{View/} & Aucun test d'interface Compose \\
\bottomrule
\end{tabularx}
\end{center}

Qualité : assertions significatives (tolérances numériques, messages explicites), aucun
\code{@Ignore}, desktopTest sur vrai SQLite en fichier temporaire. Fragilités mineures :
\code{Thread.sleep(5)} (\fichier{BackupServiceTest.kt:89}), latence aléatoire volontaire
(\fichier{PlanEvolutifInteractionTest.kt:28}), \code{GuideScreenshotDriver} ouvre la vraie base
utilisateur si une variable d'environnement manque. Les six tests iOS ne sont jamais exécutés en
CI. Les 723 tests communs tournent deux fois (desktop et Android). Aucune mesure de couverture.

\section{Build, CI et outillage}

\subsection*{Ce qui est fait correctement}

Workflow unique commenté, \code{concurrency}, permissions minimales, \code{paths-ignore} pour la
documentation, validation des wrappers, contrôle de dérive des schémas Room, iOS conditionnel,
paquet Debian conditionné aux tests verts avec contrôle AVX2, job agrégateur, Dependabot mensuel,
versions 3.3.18 et 318 cohérentes entre Android, iOS, desktop et changelog.

\subsection*{Constats}

\begin{itemize}
\item \sev{Majeur} --- Aucun lint, analyse statique (detekt, ktlint, Android Lint), ni couverture
  (Kover) ; aucun bloc \code{lint \{\}} ni \code{composeCompiler \{ reports \}}.
\item \sev{Majeur} --- Aucune publication sur tag : pas de \code{signingConfigs} Android dans Gradle,
  le \code{.deb} n'est qu'un artefact de 30 jours, signature macOS par script externe.
\item \sev{Majeur} --- \code{kotlinx-coroutines 1.7.3} (2023) avec Kotlin 2.2.20 ;
  \code{room 2.7.0-rc01} en production ; \code{jvmTarget} 11 sur Android et non fixé sur desktop
  (JDK 21 en CI) ; options \code{-X} Kotlin/Native appliquées à la cible Android.
\item \sev{Mineur} --- Catalogue Gradle encombré : une vingtaine d'alias jamais référencés
  (navigation, splitpane alpha, testng, jansi, realm, assertk, junit5, datetime 0.4.1,
  compose-plugin 1.6.11) ; dépendances déclarées sans import (paging, datastore, webview,
  reorderable, qr-kit, ui-util) ; \code{koalaplot} et \code{material-icons} codés en dur à une
  version différente du catalogue ; \code{okio} déclaré deux fois.
\item \sev{Mineur} --- \fichier{proguard-rules.pro} conserve \code{Data.**}, \code{DataBase.**},
  \code{Enumer.**}, \code{kotlin.**}, \code{io.ktor.**}, \code{okio.**} en entier : le shrinking
  est largement neutralisé ; ProGuard desktop désactivé.
\item \sev{Mineur} --- Commentaires CI périmés (« 40 + 9 fichiers de test », « 35 schémas ») ;
  cache Gradle jamais alimenté pour \code{master} et \code{revetBackUp} ; checkout LFS répété sur
  quatre jobs.
\end{itemize}

\section{Hygiène du dépôt, documentation, localisation}

\begin{itemize}
\item \sev{Critique} --- \fichier{composeApp/release/composeApp-release-000.aab} (22,8 Mo) et
  \fichier{-1-0-1.aab} (17,1 Mo) commités en blobs bruts bien que le répertoire et \code{*.aab}
  soient ignorés et suivis par LFS.
\item \sev{Critique} --- \fichier{vetnutri\_export\_init.json} (20,7 Mo) en cinq exemplaires :
  quatre identiques (commonMain, androidMain, deux dans iosApp) et une copie datée
  « 20-06-01-825 » différente et non référencée par Xcode. Neuf versions successives dans
  l'historique, soit environ 270 Mo pour les quinze plus gros blobs. Les
  \fichier{strings\_*.json} existent en trois copies déjà divergentes entre commonMain et Android.
\item \sev{Majeur} --- Jeux de données bruts à la racine : \fichier{food.json} (15 Mo),
  \fichier{animaux.json}, \fichier{ani.json}, \fichier{oza2.json}, \fichier{test\_vbnr.json},
  \fichier{docs/Untitled.json} (copie de \fichier{references.json}) ;
  \fichier{View/vetfood.json} (1,7 Mo) dans le dossier source Kotlin.
\item \sev{Majeur} --- Répertoires \fichier{commonMain/} et \fichier{iosMain/} à la racine, hors
  de tout source set (copies de KoalaPlot et d'un générateur QR) ; sept fichiers Kotlin orphelins à
  la racine dont un fichier d'un octet ; fichier \fichier{Task} vide ; journaux
  \fichier{ios\_build.log}, \fichier{ios\_build\_final.log}, \fichier{build\_log.txt} ; icônes
  en cinq copies identiques dont deux « copie » ; \fichier{composeApp/bin/} suivi ; deux paquets R
  (\fichier{vetnutriDogR}, \fichier{vetnutriExamR}) sans lien de build.
\item \sev{Mineur} --- 18 fichiers Markdown à la racine, 14 issus du même commit du 1\textsuperscript{er}
  août 2026, quatre se recouvrant sur la mémoire et les performances ; 27 autres dans
  \fichier{docs/}. \fichier{CLAUDE.md} périmé sur ses compteurs ; \fichier{.gitignore} sans
  \code{*.log}.
\item \sev{Mineur} --- Localisation : 2\,378 appels \code{translate()} mais 11 constantes de
  \code{LocalizationKeys} absentes de \fichier{strings\_fr.json} (clé brute affichée) ; six
  langues sur neuf à 617 clés sur 1\,967, dont 450 non traduites ; environ 150 \code{Text("…")} en
  français dur dans les écrans récents (multiration, troupeau, copie de ration) ; 29 fichiers de
  vue sans aucun appel de localisation.
\end{itemize}

\section{Plan d'action proposé}

Par ordre de priorité, avec l'effort estimé.

\begin{enumerate}
\item \sev{Données personnelles et binaires} (1 h + réécriture d'historique) : retirer
  \fichier{ani.json}, \fichier{animaux.json}, \fichier{food.json}, les deux AAB, la copie datée
  du JSON initial ; purger l'historique avec \code{git filter-repo} ; anonymiser tout jeu de
  données réintroduit.
\item \sev{Corrections de correction} (½ jour) : remplacer \code{NumberUtils.format} par
  \code{TextUtils.formatDecimal} ; rendre \code{ALIMENTS.refAlimUnif} nullable ou \code{RESTRICT}
  avec \code{deleteFood} en transaction et erreur affichée ; synchroniser
  \code{EquationEvaluator.expressionCache} et remplacer la clé \code{hashCode} ; unifier le poids
  métabolique sur \code{CalculMetabolique} et le tester.
\item \sev{Import et restauration} (1 jour) : \code{importAll} et \code{importFoods} dans une
  transaction Room, \code{try/catch} sur les équations, vérification de \code{envelope.version},
  erreurs d'insertion de nutriments comptées ; aligner le texte du dialogue de restauration sur le
  comportement (fusion) ou implémenter un vrai remplacement ; remonter les erreurs des neuf
  fonctions de \code{AnimalDetailViewModel} qui avalent \code{updateConsultation}.
\item \sev{Accès base} (1,5 jour) : compteurs et identifiants par \code{COUNT(*)} et
  \code{getAllFoodIds()} ; \code{getAllReferenceEv()}, poids, consultations et rations en requêtes
  \code{IN} ; sauvegarde automatique conditionnée à un indicateur de modification sans
  \code{forceRefresh()} ; suppression de l'énergie par lot.
\item \sev{Thread principal et composition} (1,5 jour) : \code{withContext(Default)} autour des
  analyses de ration et de l'indexation ; maps \code{by lazy} dans \code{Espece} ; filtre et tri de
  \code{FoodSearchComponent} dans le ViewModel avec debounce ; \code{key} sur les 48 listes ;
  ouverture Room hors \code{onCreate} sur Android.
\item \sev{Frontière vue / repository} (3 jours) : repositories privés dans les ViewModels,
  suppression de \code{AppNavRepositories}, ViewModels Conseils, Recettes, Préférences, Startup,
  Export animal ; règle métier d'arrondi et calculs graphiques déplacés dans \fichier{Data/} ;
  base commune de ViewModel avec \code{clear()} et abandon de \code{: ViewModel()} sans store.
\item \sev{Découpage des monolithes} (3 jours, par étapes) : \code{SettingsView} en sections
  (l'infrastructure existe), \code{RationsView} avec un \code{sealed class} de dialogue,
  \code{AnimalDetailView} en sections partagées entre layouts, \code{ImportUtils} en quatre
  objets, \code{AnimalDetailViewModel} en quatre ViewModels autour d'une session animal.
\item \sev{Secrets et confidentialité} (1 jour + décision) : clés jsonbin derrière un relais ou
  chiffrement obligatoire AES-GCM avec bins éphémères ; \code{allowBackup} ou règles d'exclusion
  Android ; variable d'environnement de base désactivée en production.
\item \sev{Chaîne de build} (1 jour) : coroutines 1.10.x, Room stable, \code{jvmToolchain}
  commun, retrait des options \code{-X} Kotlin/Native et des dépendances inutilisées, règles R8
  resserrées, JSON initial en une seule source copiée au build (plus de double embarquement dans
  l'APK), parsing iOS en flux.
\item \sev{CI} (1 jour) : detekt ou ktlint, Android Lint, Kover avec seuil, tests iOS exécutés,
  tests communs exécutés une seule fois, job de release sur tag avec signature par secrets.
\item \sev{Tests manquants} (2 jours) : \code{CalculMetabolique}, \code{PdfExporter},
  \code{StartupService}, \code{MultiRationViewModel}, \code{FoodEditViewModel},
  \code{DatabaseAnimalRepository}, chaîne de migrations jusqu'à 41 vérifiée explicitement, test
  de cohérence des clés de localisation, quelques tests \code{runComposeUiTest} sur la consultation
  et la ration.
\item \sev{Nettoyage et documentation} (½ jour) : code mort listé en section~\ref{sec:architecture},
  répertoires racine orphelins, journaux, icônes dupliquées ; archiver les quatre documents
  d'optimisation contredits par le code ; mettre à jour \fichier{CLAUDE.md} (compteurs, Room,
  règle \code{withContext}) et \fichier{GUIDE\_COMPATIBILITE\_16KB.md}.
\item \sev{Navigation} (2 jours, optionnel) : pile \code{ArrayDeque<Screen>} avec écrans
  paramétrés, suppression de \code{detailReturnScreen} et \code{TestYellowBox}.
\end{enumerate}

Les quatre premiers points tiennent en moins d'une semaine et traitent tout ce qui touche
l'exactitude des résultats affichés, l'intégrité des données et la confidentialité. Les points
suivants sont des investissements de maintenabilité qui conditionnent le coût de chaque évolution
future.

\end{document}
